% ============================================================
% §04 — Domain Kit Model
% Author: Ken (Software Architect)
% Integrated by: Leslie (LaTeX Engineer)
% Last updated: 2026-04-20
% ============================================================

\chapter{Domain Kit Model}
\label{ch:domain-kits}

% ============================================================
\section{Purpose and Scope}
\label{sec:kit-purpose}
% ============================================================

The gert core runtime is deliberately domain-agnostic.  It executes steps,
enforces governance, captures evidence, and writes traces---but it carries no
opinion about what kind of work a runbook describes.  This is a feature, not an
omission: a domain-agnostic kernel is the prerequisite for a stable, long-lived
runtime contract.

Domain Kits are the mechanism by which domain-specific semantics enter the
system without contaminating the kernel.  A Domain Kit is a higher-level
semantic layer built atop core runtime primitives.  It provides five things:

\begin{enumerate}
  \item \textbf{Authoring schemas.}  A Kit defines a domain-specific YAML
    vocabulary---specialised step types, field structures, validation
    constraints---that is natural for authors working in that domain.  An
    workflow Kit might offer a \texttt{select-priority} step type with severity
    and escalation fields; a compliance Kit might offer a
    \texttt{approval-request} type with approver-matrix fields.

  \item \textbf{Validation rules.}  A Kit contributes domain-specific semantic
    validators that run at parse time or plan time.  These validators enforce
    invariants that the core schema cannot express: ``a cleanup step must
    follow every provision step,'' ``every runbook must declare an owner,''
    ``escalation timeout must not exceed configured limit.''

  \item \textbf{Lowering (compilation into core primitives).}  Kit-specific
    step types do not reach the runtime.  A Kit compiler transforms (lowers)
    its domain-specific AST nodes into sequences of core step types
    (\texttt{cli}, \texttt{tool}, \texttt{manual}, \texttt{branch},
    \texttt{iterate}) before the planner produces an ExecutionPlan.  The
    runtime executes only core primitives; it is never aware that a Kit was
    involved.

  \item \textbf{Projections and read models.}  A Kit may define read-only
    projections over run traces.  A compliance Kit might project an ``audit
    log'' view from the raw trace events; a workflow Kit might project a
    ``timeline'' view.  Projections are computed after execution; they do not
    affect runtime behaviour.

  \item \textbf{Testing contracts.}  A Kit defines what ``correct'' means for
    its domain: fixture runbooks, expected lowered output, expected validation
    errors, expected projection shapes.  These contracts are the Kit's
    acceptance test suite.
\end{enumerate}

Domain Kits do not contribute runtime capabilities.  They do not register
tools, subscribe to events, or modify execution behaviour.  Those are extension
concerns (\S\ref{ch:extension}).  A Kit operates entirely at the authoring and
analysis layer---it shapes what authors write and what reviewers see, not what
the engine does.

% ============================================================
\section{Architectural Narrative}
\label{sec:kit-narrative}
% ============================================================

The architecture draws a hard boundary between three concerns: what the
engine executes (core runtime), what capabilities are available at runtime
(extensions and tools), and what vocabulary authors use to express intent
(Domain Kits).  The first two are well-defined in the existing design; the
third has been implicit until now.

Without Domain Kits, domain semantics have two bad places to live.  They can
leak into the core schema, which makes the kernel unstable---every new domain
(compliance audits, workflow automation, data migration, system provisioning)
adds fields and step types to the core, bloating it into an ever-expanding
union type.  Or they can be pushed into extensions, which conflates authoring
concerns with runtime concerns---an extension that contributes domain-specific
validators, authoring schemas, and projections is doing fundamentally different
work than an extension that contributes tools or event subscriptions.

Domain Kits solve this by providing a clean, well-defined home for domain
semantics.  The design principle is:

\begin{description}
  \item[Core owns execution.] The core runtime
    (\S\ref{ch:architecture}), governance layer (\S\ref{ch:governance}),
    evidence/tracing system (\S\ref{ch:evidence}), and event model
    (\S\ref{ch:events}) are stable, domain-agnostic infrastructure.

  \item[Kits own vocabulary.] Authoring schemas, domain validation rules,
    lowering/compilation, projections, and testing contracts belong to Kits.  A
    Kit is the right place to encode ``what does a well-formed workflow runbook
    look like?'' or ``what fields does an approval-request step require?''

  \item[Extensions own capabilities.] Tool registration, event subscription,
    schema extension (namespaced \texttt{x-*} fields), policy contribution, and
    provider registration belong to extensions (\S\ref{ch:extension}).

  \item[Tools own effects.] Executable actions---invoking binaries, calling
    APIs, performing side-effecting work---belong to the tool runtime
    (\S\ref{ch:tools}).
\end{description}

This separation preserves the kernel's stability guarantee: the core schema and
runtime contract do not change when a new domain is added.  The Kit compiles
domain-specific vocabulary down to core primitives; the runtime never sees
Kit-specific nodes.

% ============================================================
\section{The Clean Kernel Principle}
\label{sec:kit-clean-kernel}
% ============================================================

The key architectural invariant is: \textbf{the runtime never executes
Kit-specific step types.}  All domain-specific authoring nodes are lowered
(compiled) into core primitives before the planner sees them.

The compilation pipeline is:

\begin{verbatim}
Kit YAML → Kit Parser → Kit AST → Kit Compiler → Core YAML → Core Parser → ExecutionPlan
\end{verbatim}

Or equivalently, when integrated into the gert pipeline:

\begin{enumerate}
  \item The user authors a runbook using Kit vocabulary (e.g.,
    \texttt{apiVersion: runbook/v1+ops/v1}).
  \item The Kit's parser validates domain-specific fields and produces a
    Kit-level AST.
  \item The Kit's compiler lowers the AST into a standard \texttt{runbook/v1}
    document using only core step types.
  \item The core parser and planner process the lowered document normally.
  \item The runtime executes the plan with no knowledge that a Kit was
    involved.
\end{enumerate}

This is a compile-time transformation, not a runtime interception.  The lowered
output is a valid \texttt{runbook/v1} document that can be inspected,
version-controlled, and executed without the Kit installed.  This property is
critical for portability: a Kit-authored runbook can always be reduced to its
core equivalent.

The Kit compiler must be deterministic: the same Kit input must always produce
the same core output.  This ensures that \texttt{gert test} scenarios written
against lowered output remain stable across Kit versions.

% ============================================================
\section{Layer Relationships}
\label{sec:kit-layers}
% ============================================================

The system has four distinct layers.  Each layer has a defined
responsibility boundary and a defined relationship to the others.

\begin{table}[ht]
\centering
\caption{Four-layer responsibility matrix}
\label{tab:kit-layers}
\small
\begin{tabular}{p{0.15\linewidth}p{0.30\linewidth}p{0.25\linewidth}p{0.18\linewidth}}
\toprule
\textbf{Layer} & \textbf{Owns} & \textbf{Does NOT own} & \textbf{Runs at} \\
\midrule
Core Runtime &
  Execution, state machine, governance enforcement, evidence, event emission, trace writing &
  Domain vocabulary, tool implementations, UI rendering &
  Run time \\
\midrule
Domain Kits &
  Authoring schemas, domain validation, lowering/compilation, projections, testing contracts &
  Execution, runtime capabilities, side effects &
  Author time / compile time \\
\midrule
Extension Runtime &
  Runtime capabilities: tool registration, event subscription, schema extension (\texttt{x-*}), policy contribution, provider registration &
  Authoring vocabulary, domain semantics, execution control &
  Run time (out-of-process) \\
\midrule
Tool Runtime &
  Side-effecting actions: binary execution, API calls, MCP tool invocation &
  Governance, evidence, schema, authoring vocabulary &
  Run time (per-invocation or persistent process) \\
\bottomrule
\end{tabular}
\end{table}

Key distinctions:

\begin{description}
  \item[Domain Kits are NOT extensions.] Extensions contribute runtime
    capabilities via the capability grant model
    (\S\ref{sec:ext-capabilities}).  Kits contribute authoring semantics via
    compilation.  An extension runs as a child process during execution; a Kit
    runs as a compiler before execution.  They serve fundamentally different
    purposes and operate at different phases.

  \item[Domain Kits are NOT tools.] Tools perform side-effecting work.  Kits
    are pure transformations with no side effects.  A Kit compiler reads YAML
    and writes YAML; it never invokes binaries, calls APIs, or modifies system
    state.

  \item[Extensions may accompany a Kit but are distinct.] A Kit distribution
    (e.g., a ``compliance workflow pack'') may ship both a Domain Kit (authoring
    schemas, lowering rules) and an extension (workflow-specific tools,
    PagerDuty providers).  These are bundled for convenience but are
    architecturally separate: the Kit operates at compile time, the extension
    operates at run time.
\end{description}

% ============================================================
\section{Kit Manifest and Compatibility}
\label{sec:kit-manifest}
% ============================================================

A Domain Kit ships a \texttt{gert-kit.yaml} manifest alongside its compiler.
The manifest declares identity, compatibility, and the Kit's schema
contributions.

\begin{minted}{yaml}
# gert-kit.yaml — Domain Kit Manifest
apiVersion: kit/v1

meta:
  name: gert.workflow                # Reverse-DNS namespaced identifier
  version: "1.0.0"                   # Semver
  description: "Workflow automation authoring kit"
  author: "Gert Project <gert@ormasoft.dev>"

compatibility:
  gert-core: ">=2.0.0 <3.0.0"       # Minimum gert core version (semver range)

schema:
  apiVersion-suffix: "workflow/v1"   # Runbooks authored with this Kit use
                                     # apiVersion: runbook/v1+workflow/v1
  json-schema: "./schemas/workflow-v1.json"  # Kit-specific JSON Schema overlay

compiler:
  executable: "./bin/gert-kit-workflow"   # Kit compiler binary
  args: ["compile"]
  input-format: yaml                 # Kit YAML in, core YAML out
  output-format: yaml

step-types:                          # Domain-specific step types this Kit defines
  - name: approval-request
    lowers-to: [manual, tool]        # Core types this compiles into
  - name: cleanup
    lowers-to: [cli, branch]
  - name: select-priority
    lowers-to: [manual, branch]

validators:                          # Domain-specific validation rules
  - name: workflow/cleanup-follows-provision
    phase: plan                      # parse | plan
    description: "Every provision step must be followed by a cleanup step"
  - name: workflow/owner-required
    phase: parse
    description: "Every workflow runbook must declare an owner in meta.roles"

projections:                         # Read models over trace data
  - name: workflow/audit-log
    description: "Workflow timeline derived from trace events"
  - name: workflow/evidence-summary
    description: "Audit-ready evidence summary grouped by governance primitive"
\end{minted}

\paragraph{Compatibility and versioning rules.}

\begin{itemize}
  \item A Kit declares a semver range for the minimum gert core version in
    \texttt{compatibility.gert-core}.  If the host version falls outside this
    range, the Kit is rejected at compile time with a structured error.

  \item Breaking changes to a Kit's authoring schema require a Kit major
    version bump.  The \texttt{apiVersion-suffix} encodes the Kit's major
    version (e.g., \texttt{workflow/v1} $\to$ \texttt{workflow/v2 (later major)}).  This is
    independent of the gert core version.

  \item Additive changes (new optional fields, new step types) within a Kit
    major version are backward-compatible and do not require a version bump to
    the suffix.

  \item The Kit's \texttt{json-schema} artifact carries a semver artifact
    version for minor/patch tracking, following the same pattern as core schema
    artifacts (\S\ref{subsec:versioning-strategy}).
\end{itemize}

\paragraph{Local-first constraints.}

Domain Kits are portable semantic layers.  They run on any supported gert
host---a local laptop, a CI runner, a self-hosted server---without requiring
network access, cloud services, or central registries.  A Kit is a directory
containing a manifest, a compiler binary, and schema artifacts.  It can be
vendored into a repository, distributed as a tarball, or installed from a
package manager.  There is no Kit registry service, no Kit telemetry, and no
phone-home behaviour.  This aligns with gert's local-first design philosophy
(\S\ref{ch:overview}).

% ============================================================
\section{Extension Runtime Audit}
\label{sec:kit-ext-audit}
% ============================================================

The following table audits concerns currently described in the extension
runtime chapter (\S\ref{ch:extension}) and evaluates whether each is properly
an extension concern, a Domain Kit concern, or shared.

\begin{longtable}{p{0.15\linewidth}p{0.13\linewidth}p{0.14\linewidth}p{0.46\linewidth}}
\toprule
\textbf{Concern} & \textbf{Current location} & \textbf{Recommendation} & \textbf{Rationale} \\
\midrule
\endfirsthead

\toprule
\textbf{Concern} & \textbf{Current location} & \textbf{Recommendation} & \textbf{Rationale} \\
\midrule
\endhead

\midrule
\multicolumn{4}{r}{\footnotesize\itshape Continued on next page} \\
\bottomrule
\endfoot

\bottomrule
\endlastfoot

Domain-specific validators &
  Implicit in \texttt{capability/\allowbreak schema-extension} &
  \textbf{Move to Kit.} &
  Validation rules that enforce domain invariants (e.g., ``cleanup must follow
  provision'') are authoring-time concerns.  They belong in the Kit's validation
  phase, not in a runtime extension.  The extension
  \texttt{capability/\allowbreak schema-extension} should be reserved for namespaced
  field contributions (\texttt{x-*}), not for semantic validation of domain
  workflows. \\

Authoring schema registration &
  \texttt{capability/\allowbreak schema-extension} &
  \textbf{Clarify boundary.} &
  Extensions contribute namespaced runtime-observable fields (\texttt{x-acme.*}).
  Kits contribute domain-specific authoring vocabulary (new step types, required
  field structures).  These are different operations.  Add a clarifying note to
  \S\ref{sec:ext-capabilities} that \texttt{capability/\allowbreak schema-extension} is
  for runtime-observable namespace extensions, not for authoring-vocabulary
  contributions. \\

Projections / read models &
  Not currently addressed &
  \textbf{Kit concern.} &
  Projections are pure, read-only transformations over trace data.  They are
  computed after execution and do not require runtime capabilities.  They belong
  in the Kit layer.  Extensions that need to observe execution should use
  \texttt{capability/\allowbreak event-subscription}; post-hoc analysis belongs to
  Kits. \\

Compilation hooks &
  Not currently addressed &
  \textbf{Kit concern.} &
  The lowering/compilation step is a pure transformation that happens before the
  runtime starts.  It has no runtime footprint and requires no capabilities.
  Compilation is entirely within the Kit layer. \\

Domain testing fixtures &
  Not currently addressed &
  \textbf{Kit concern.} &
  Test fixtures for domain-specific authoring (example runbooks, expected lowered
  output, expected validation errors) are Kit deliverables.  They do not interact
  with the extension runtime. \\

Tool registration &
  \texttt{capability/\allowbreak tool-registration} &
  \textbf{Stays in extension runtime.} &
  Tool registration is a runtime capability.  A Kit distribution may ship an
  accompanying extension that registers domain-specific tools, but tool
  registration itself remains an extension concern. \\

Policy contribution &
  \texttt{capability/\allowbreak policy-contribution} &
  \textbf{Stays in extension runtime.} &
  Policy rules that are evaluated at run time (pre-step, pre-run) are runtime
  concerns.  Domain-specific validation rules that are evaluated at parse/plan
  time belong to Kits; runtime policy enforcement belongs to extensions. \\

Event subscription &
  \texttt{capability/\allowbreak event-subscription} &
  \textbf{Stays in extension runtime.} &
  Event observation is inherently a runtime activity.  No change needed. \\

\end{longtable}

\paragraph{Summary of changes needed in existing sections.}

\begin{enumerate}
  \item Add a clarifying paragraph to \texttt{capability/\allowbreak schema-extension}
    (\S\ref{sec:ext-capabilities}) noting the boundary: extensions contribute
    namespaced runtime fields; Kits contribute authoring vocabulary and domain
    step types.

  \item Add a forward reference from the extension runtime chapter to the
    Domain Kit chapter for readers who arrive at the extension spec looking for
    domain-specific validation or projections.

  \item No capabilities need to be moved or removed from the extension
    taxonomy.  The taxonomy is correct for runtime concerns; the issue was that
    domain/authoring concerns had no defined home.  Domain Kits provide that
    home.
\end{enumerate}

% ============================================================
\section{Domain Kit Examples}
\label{sec:kit-examples}
% ============================================================

The Domain Kit model enables multiple domain-specific vocabularies to coexist
without contaminating the gert core. Below are representative examples of how
different domains might structure their kits.

\paragraph{Reference implementations.}
The gert project maintains several reference Domain Kits to demonstrate the
model in practice. These are NOT part of gert core, but are distributed
separately as example implementations:

\begin{description}
  \item[\texttt{gert.workflow}] --- General workflow automation kit with
    approval-request, cleanup, and select-priority step types. Demonstrates
    basic Kit compiler structure.
  
  \item[\texttt{gert.compliance}] --- Compliance and audit-focused kit with
    evidence collection, audit-log projections, and policy enforcement patterns.
    
  \item[\texttt{gert.migration}] --- Data migration kit with rollback semantics,
    checkpoint/restart, and validation step types.
    
  \item[\texttt{gert.ops}] --- Operations runbook kit implementing the DRI
    (Directly Responsible Individual) model for incident response and change
    management. See the \emph{DRI Domain Kit Manual} for details.
\end{description}

The architectural principle remains firm: \textbf{gert core contains zero
domain-specific vocabulary}. All step types in the core schema (\texttt{cli},
\texttt{manual}, \texttt{tool}, \texttt{branch}, \texttt{iterate}) are
domain-agnostic execution primitives. Domain semantics live exclusively in
separately-distributed Kits.

For detailed Kit implementation guides, see the \emph{Domain Kit Development
Guide} (separate document).

% ============================================================
\section{Kit Composition and Layering}
\label{sec:kit-composition}
% ============================================================

A specialised Domain Kit may depend on a foundational Domain Kit.  The
composition mechanism is a shared \texttt{CompilerRegistry} that each kit
populates at process start-up.  The output invariant---a flat gert core
\texttt{ExecutionPlan}---is never broken regardless of how many Kit layers are
involved, because every step compiler at every layer ultimately emits
\texttt{[]schema.FlowNode} using only core primitive types.

% ------------------------------------------------------------
\subsection{The Composition Model}
\label{subsec:kit-composition-model}
% ------------------------------------------------------------

\subsubsection{What a Kit Exposes}

A foundational kit exposes three integration surfaces:

\begin{enumerate}
  \item \textbf{Model types.}  Domain value objects (\texttt{Task},
    \texttt{Person}, \texttt{Approval}, \texttt{Chore}) as plain Go structs
    with YAML tags and validation constraints.  These are the ``nouns'' of the
    domain vocabulary.

  \item \textbf{Step compiler functions.}  One \texttt{StepCompilerFunc}
    implementation per step type.  A step compiler accepts a raw
    \texttt{*yaml.Node} (the step body as authored) and returns
    \texttt{([]schema.FlowNode, error)}.  The returned nodes are core gert
    primitives only.

  \item \textbf{A \texttt{KitBundle}.}  A named collection of
    (step-type~$\to$~compiler-function) pairs ready to be merged into a
    \texttt{CompilerRegistry}.  This is the integration surface for downstream
    kits.
\end{enumerate}

The shared infrastructure types live in
\texttt{gert/pkg/kitruntime}:

\begin{minted}{go}
// StepCompilerFunc compiles one Kit-level step into core FlowNodes.
type StepCompilerFunc func(
    ctx   context.Context,
    node  *yaml.Node,
    scope CompileScope,
) ([]schema.FlowNode, error)

// CompileScope carries read-only context available during compilation.
type CompileScope struct {
    Vars       map[string]string
    Governance *schema.GovernanceConfig
    KitMeta    KitMeta
    Registry   *CompilerRegistry
}

// KitBundle is what a foundational kit exports.
type KitBundle struct {
    Name      string
    Version   string
    Compilers map[string]StepCompilerFunc // key: qualified type, e.g. "household.chore"
}

// CompilerRegistry holds all registered step compilers, keyed by qualified type.
type CompilerRegistry struct {
    compilers map[string]StepCompilerFunc
    owners    map[string]string // step key -> kit name, for diagnostics
}

func NewRegistry() *CompilerRegistry { ... }

// Merge adds all compilers from a KitBundle into the registry.
// Panics on duplicate key — collision is a programming error.
func (r *CompilerRegistry) Merge(b KitBundle) { ... }

// Dispatch looks up and calls the compiler for a given qualified step type.
func (r *CompilerRegistry) Dispatch(
    ctx      context.Context,
    stepType string,
    node     *yaml.Node,
    scope    CompileScope,
) ([]schema.FlowNode, error) { ... }
\end{minted}

\subsubsection{How the Specialised Kit References the Foundational Kit}

The specialised kit (\texttt{gert-kit-home}) imports the foundational kit
(\texttt{gert-kit-household}) as a normal Go module dependency:

\begin{minted}{go}
// gert-kit-home/go.mod
module github.com/ormasoftchile/gert-kit-home

require (
    github.com/ormasoftchile/gert-kit-household v0.3.0
    github.com/ormasoftchile/gert/v1             .0
)
\end{minted}

At start-up the specialised kit builds its registry by merging the foundational
kit's bundle and then adding its own compilers:

\begin{minted}{go}
// gert-kit-home/pkg/compiler/registry.go
func BuildRegistry() *kitruntime.CompilerRegistry {
    r := kitruntime.NewRegistry()
    r.Merge(household.Bundle()) // foundational kit contributes its compilers
    r.Merge(homeBundle())       // specialised kit adds its own
    return r
}
\end{minted}

\texttt{household.Bundle()} is the exported integration surface of
\texttt{gert-kit-household}:

\begin{minted}{go}
// gert-kit-household/pkg/compiler/bundle.go
func Bundle() kitruntime.KitBundle {
    return kitruntime.KitBundle{
        Name:    "household",
        Version: "0.3.0",
        Compilers: map[string]kitruntime.StepCompilerFunc{
            "household.chore":   compileChore,
            "household.approve": compileApprove,
            "household.notify":  compileNotify,
        },
    }
}
\end{minted}

\subsubsection{Composition Phases}

\begin{description}
  \item[Go module dependency — build time.]  The type system enforces the
    contract.  A signature change in \texttt{household.Bundle()} is a compile
    error in \texttt{gert-kit-home}.

  \item[Registry merge --- load time.]  When the kit compiler process starts
    (or when a test calls \texttt{BuildRegistry()}), the registry is assembled.
    Step type dispatch is resolved at this point.

  \item[Step lowering --- authoring-compile time.]  When the kit compiler
    processes a \texttt{.home.yaml} input, it runs step compilers from the
    registry.  This ``compile time'' precedes gert runtime execution.
\end{description}

% ------------------------------------------------------------
\subsection{Step Type Namespacing}
\label{subsec:kit-step-namespacing}
% ------------------------------------------------------------

\textbf{Decision: prefixed step types.}  Format: \texttt{<kit-name>.<step-type>}.

\begin{minted}{yaml}
- type: household.chore
  ...
- type: household.approve
  ...
- type: home.morning-routine
  ...
\end{minted}

Two alternatives were considered and rejected:

\begin{description}
  \item[Implicit disjoint sets (rejected).]  Bare type names rely on kit authors
    never colliding.  With two kits in scope, \texttt{approve} could silently
    mean either \texttt{household.approve} or \texttt{core.approve}.  The
    \texttt{Merge()} panic catches exact key matches, but the convention breaks
    down in a public kit ecosystem.

  \item[\texttt{kit:} field (rejected).]  Splitting identity across two YAML
    fields (\texttt{kit: household}, \texttt{type: chore}) makes single-field
    lookups ambiguous.  The prefix form \texttt{household.chore} is
    self-contained and survives copy-paste.
\end{description}

The prefix form wins on four grounds: (1) a single \texttt{type:} field carries
full identity; (2) the registry key is the YAML \texttt{type} value, so no
translation is needed at dispatch; (3) domain context is visible in every
runbook listing; (4) collisions are caught at start-up via \texttt{Merge()}.

\paragraph{Core step types are unqualified.}  The gert core step types
(\texttt{cli}, \texttt{manual}, \texttt{tool}, \texttt{approve},
\texttt{branch}, \texttt{iterate}) are not managed by any Kit and remain
unqualified.  The rule is unambiguous: unqualified~= core primitive;
qualified~= Kit-authored.

% ------------------------------------------------------------
\subsection{Compiler Delegation}
\label{subsec:kit-compiler-delegation}
% ------------------------------------------------------------

\textbf{Decision: Option~B --- shared \texttt{CompilerRegistry} with dispatch
by step type.}

\paragraph{Option A: direct function calls (rejected).}

The specialised kit would enumerate every foundational step type it might
encounter:

\begin{minted}{go}
// Option A (rejected)
switch step.Type {
case "household.chore":
    nodes, err = household.CompileChore(ctx, node, scope)
case "household.approve":
    nodes, err = household.CompileApprove(ctx, node, scope)
// ...must list every household type explicitly
\end{minted}

This tightly couples the specialised kit to the foundational kit's internal
type inventory.  Any new step type added to \texttt{household} requires a
corresponding \texttt{case} in \texttt{home}.

\paragraph{Option C: \texttt{BaseCompiler} embedding (rejected).}

\begin{minted}{go}
// Option C (rejected)
type HomeCompiler struct {
    household.BaseCompiler // embedding
}
\end{minted}

Go embedding works for single inheritance but breaks down when a specialised
kit depends on multiple foundational kits.  It also makes the foundational
compiler hard to replace with a test double.

\paragraph{Option B: shared registry dispatch (adopted).}

The \texttt{CompilerRegistry} is a simple map with a dispatch function.  Each
kit contributes its entries at start-up.  The specialised kit's compiler calls
\texttt{registry.Dispatch(stepType,~...)} for every step, regardless of which
kit owns it:

\begin{minted}{go}
// Option B
func (c *HomeCompiler) CompileStep(
    ctx   context.Context,
    step  RawStep,
    scope Scope,
) ([]schema.FlowNode, error) {
    return c.registry.Dispatch(ctx, step.Type, step.Body, scope)
}
\end{minted}

New step types added to the foundational kit are automatically available: the
specialised kit calls \texttt{r.Merge(household.Bundle())} and receives them
all.

\textit{Trade-off acknowledged:} all bundles must be merged before the first
\texttt{Dispatch()} call.  This is a start-up invariant enforced by making
\texttt{BuildRegistry()} a required, non-lazy setup step.

% ------------------------------------------------------------
\subsection{Model Composition}
\label{subsec:kit-model-composition}
% ------------------------------------------------------------

\textbf{Decision: Go embedding for model types, plus direct import of
foundational types.}

Kit model types are \emph{value objects}---data carriers, not behaviour objects.
Go embedding is the idiomatic way to extend value types without losing type
safety.  Interfaces add indirection without benefit at the model layer; they
are reserved for the compiler functions (\texttt{StepCompilerFunc}), not the
data types.

\begin{minted}{go}
// gert-kit-household/pkg/model/types.go
type Chore struct {
    ID          string  `yaml:"id"`
    Label       string  `yaml:"label"`
    Description string  `yaml:"description"`
    Zone        string  `yaml:"zone"`
    Frequency   Cadence `yaml:"cadence"`
    Assignee    string  `yaml:"assignee,omitempty"`
}

// gert-kit-home/pkg/model/types.go
import household "github.com/ormasoftchile/gert-kit-household/pkg/model"

type MorningRoutine struct {
    ID    string        `yaml:"id"`
    Label string        `yaml:"label"`
    Steps []RoutineStep `yaml:"steps"`
}

type RoutineStep struct {
    Type  string           `yaml:"type"`  // "chore" | "approve" | "notify"
    Chore *household.Chore `yaml:"chore,omitempty"`
}
\end{minted}

% ------------------------------------------------------------
\subsection{The ExecutionPlan Invariant}
\label{subsec:kit-invariant}
% ------------------------------------------------------------

\textbf{Invariant:} regardless of how many Kit layers are involved, the output
of the full compilation chain is always a flat gert core
\texttt{ExecutionPlan}.  Kit-specific step types never reach the gert runtime.

The invariant is enforced at three levels:

\begin{enumerate}
  \item \textbf{Type signature.}  \texttt{StepCompilerFunc} returns
    \texttt{[]schema.FlowNode}.  \texttt{schema.FlowNode} is the core union
    type; it has no Kit-level variant.  The type system prevents Kit-level
    nodes from being returned.

  \item \textbf{Composite expansion.}  When \texttt{home.morning-routine}
    compiles, it does not emit \texttt{household.chore} as a
    \texttt{FlowNode}.  It calls
    \texttt{registry.Dispatch("household.chore",~...)} for each sub-step and
    concatenates the resulting core nodes:

\begin{minted}{go}
func compileMorningRoutine(
    ctx   context.Context,
    node  *yaml.Node,
    scope kitruntime.CompileScope,
) ([]schema.FlowNode, error) {
    var routine model.MorningRoutine
    if err := node.Decode(&routine); err != nil {
        return nil, err
    }
    var allNodes []schema.FlowNode
    for _, step := range routine.Steps {
        nodes, err := scope.Registry.Dispatch(
            ctx, step.QualifiedType(), step.Body, scope,
        )
        if err != nil {
            return nil, fmt.Errorf(
                "home.morning-routine step %s: %w", step.ID, err,
            )
        }
        allNodes = append(allNodes, nodes...)
    }
    return allNodes, nil // flat list of core nodes only
}
\end{minted}

  \item \textbf{Core planner boundary.}  The gert core planner accepts only
    \texttt{runbook/v1} YAML.  Kit YAML is a pre-processing artefact that
    never enters the planner.
\end{enumerate}

% ------------------------------------------------------------
\subsection{Worked Example: Household and Home Kits}
\label{subsec:kit-worked-example}
% ------------------------------------------------------------

\textbf{Kits in scope.}
\texttt{gert-kit-household} (foundational) defines \texttt{household.chore},
\texttt{household.approve}, and \texttt{household.notify}.
\texttt{gert-kit-home} (specialised) imports \texttt{household} and adds
\texttt{home.morning-routine} as a composite step type.

\subsubsection{Home Runbook Source (Kit YAML)}

\begin{minted}{yaml}
# casa-santiago.home.yaml
apiVersion: runbook/v1+home/v1
id: morning-routine-weekday
name: Weekday Morning Routine
kind: reference

vars:
  day: monday

flow:
  - type: home.morning-routine
    id: morning_core
    label: Core Morning Tasks
    steps:
      - type: household.chore
        id: make_coffee
        label: Make coffee
        zone: kitchen
        cadence: { frequency: daily }

      - type: household.chore
        id: feed_pets
        label: Feed cats
        zone: utility
        cadence: { frequency: daily }

      - type: household.approve
        id: morning_sign_off
        label: Morning checklist sign-off
        roles: [resident]
        required: 1
        timeout: 30m

      - type: household.notify
        id: done_notify
        label: Notify family
        channel: push
        message: "Morning routine complete"
\end{minted}

\subsubsection{Compiled Output (Core \texttt{runbook/v1} YAML)}

After compilation by \texttt{gert-kit-home}, no Kit types remain:

\begin{minted}{yaml}
# Compiled output — no kit types remain
apiVersion: runbook/v1
id: morning-routine-weekday
name: Weekday Morning Routine
kind: reference

flow:
  - type: manual              # household.chore -> manual step
    id: make_coffee
    title: Make coffee
    subtitle: "[kitchen] daily"

  - type: manual
    id: feed_pets
    title: Feed cats
    subtitle: "[utility] daily"

  - type: approve             # household.approve -> approve gate
    id: morning_sign_off
    title: Morning checklist sign-off
    roles: [resident]
    required: 1
    timeout: 30m

  - type: tool                # household.notify -> tool step
    id: done_notify
    title: Notify family
    tool: notify
    action: push
    args:
      message: Morning routine complete
\end{minted}

\subsubsection{Go Interface Sketch}

\begin{minted}{go}
// --- gert-kit-household/pkg/compiler/bundle.go ---

func Bundle() kitruntime.KitBundle {
    return kitruntime.KitBundle{
        Name:    "household",
        Version: "0.3.0",
        Compilers: map[string]kitruntime.StepCompilerFunc{
            "household.chore":   compileChore,
            "household.approve": compileApprove,
            "household.notify":  compileNotify,
        },
    }
}

func compileChore(
    ctx   context.Context,
    node  *yaml.Node,
    scope kitruntime.CompileScope,
) ([]schema.FlowNode, error) {
    var chore model.Chore
    if err := node.Decode(&chore); err != nil {
        return nil, fmt.Errorf("household.chore: %w", err)
    }
    step := schema.Step{
        ID:       chore.ID,
        Type:     schema.StepTypeManual,
        Title:    chore.Label,
        Subtitle: fmt.Sprintf("[%s] %s", chore.Zone, chore.Frequency.Human()),
    }
    return []schema.FlowNode{{Step: &step}}, nil
}

// compileApprove and compileNotify follow the same pattern.


// --- gert-kit-home/pkg/compiler/registry.go ---

func BuildRegistry() *kitruntime.CompilerRegistry {
    r := kitruntime.NewRegistry()
    r.Merge(household.Bundle()) // foundational kit
    r.Merge(homeBundle())       // specialised kit
    return r
}

func homeBundle() kitruntime.KitBundle {
    return kitruntime.KitBundle{
        Name:    "home",
        Version: "0.1.0",
        Compilers: map[string]kitruntime.StepCompilerFunc{
            "home.morning-routine": compileMorningRoutine,
        },
    }
}


// --- gert-kit-home/pkg/compiler/compiler.go ---

type Compiler struct {
    registry *kitruntime.CompilerRegistry
}

func New() *Compiler {
    return &Compiler{registry: BuildRegistry()}
}

// Compile transforms a .home.yaml document into a core runbook/v1 document.
func (c *Compiler) Compile(ctx context.Context, input []byte) ([]byte, error) {
    // 1. Parse input as Home Kit AST.
    // 2. For each flow step, call c.registry.Dispatch(stepType, node, scope).
    // 3. Assemble schema.Runbook with the returned []schema.FlowNode.
    // 4. Marshal to YAML and return (valid runbook/v1 YAML).
}
\end{minted}

% ------------------------------------------------------------
\subsection{Border Cases and Error Handling}
\label{subsec:kit-border-cases}
% ------------------------------------------------------------

\begin{table}[ht]
\centering
\caption{Kit composition border cases}
\label{tab:kit-border-cases}
\small
\begin{tabular}{p{0.20\linewidth}p{0.25\linewidth}p{0.45\linewidth}}
\toprule
\textbf{Scenario} & \textbf{Specified behaviour} & \textbf{Invariant / guard} \\
\midrule
Unknown step type at dispatch &
  \texttt{Dispatch()} returns a structured error; never panics. &
  Error message includes the unknown type name and the registry size:
  \texttt{kitruntime: unknown step type "X" (registry contains N types)}.
  Count aids diagnosis when the registry was built with the wrong bundles. \\
\midrule
Version conflict between kits &
  Go MVS resolves the module conflict before the binary is built---exactly
  one \texttt{Bundle()} exists in the final binary.  If the same key is merged
  twice, \texttt{Merge()} panics. &
  \textit{Merge ownership invariant:} a foundational kit exports
  \texttt{Bundle()} but never calls \texttt{r.Merge()} itself.  Only the
  top-level \texttt{BuildRegistry()} site calls \texttt{Merge()}, ensuring
  each bundle is merged exactly once. \\
\midrule
Overlapping prefixes / namespace collision &
  An exact key collision causes a start-up panic (programming error).
  Substring overlap is not detected at runtime; it is a documentation problem. &
  \textit{Prefix reservation invariant:} each kit owns its prefix exclusively.
  The \texttt{gert-kit.yaml} manifest must declare a \texttt{prefix:} field.
  The \texttt{Merge()} panic names both the existing owner and the incoming
  bundle for immediate diagnosis. \\
\midrule
Partial registration / startup order &
  If a bundle is not merged before \texttt{Dispatch()} is called, the
  ``unknown step type'' error from \S\ref{subsec:kit-border-cases} propagates
  with no silent fallback. &
  \textit{Startup order invariant:} \texttt{BuildRegistry()} must complete and
  return a fully-populated registry before any \texttt{Dispatch()} call is
  possible.  The constructor pattern (\texttt{New()~=~\&Compiler\{registry:
  BuildRegistry()\}}) enforces this. \\
\midrule
Nil or empty step body &
  The kit's input loader rejects null step bodies before calling
  \texttt{Dispatch()}.  \texttt{StepCompilerFunc} implementations may assume
  \texttt{node} is non-nil. &
  \texttt{Dispatch()} passes \texttt{node} to the compiler unchanged; it does
  not inject a synthetic node.  The loader is the validation boundary: \\
  & &
  \texttt{if node == nil \{\ return nil,\ fmt.Errorf("step \%q:\ body required but null",\ id)\ \}} \\
\bottomrule
\end{tabular}
\end{table}

% ============================================================
\section{Non-Goals}
\label{sec:kit-nongoals}
% ============================================================

Domain Kits are a precisely scoped mechanism.  The following are explicitly not
in scope:

\begin{description}
  \item[Kits are not extensions.] A Kit does not contribute runtime
    capabilities.  It does not register tools, subscribe to events, modify
    execution behaviour, or run as a child process during execution.  The
    extension capability taxonomy (\S\ref{sec:ext-capabilities}) does not
    apply to Kits.

  \item[Kits are not tools.] A Kit does not perform side-effecting work.  It
    does not invoke binaries, call APIs, or modify system state.  The Kit
    compiler is a pure function: YAML in, YAML out.

  \item[Kits are not plugins in a generic sense.] Gert does not have a
    ``plugin system.''  It has three distinct extension points---Domain Kits
    (authoring semantics), extensions (runtime capabilities), and tools
    (effectful actions)---each with a different contract, lifecycle, and trust
    model.  Calling any of these ``plugins'' obscures the architectural
    boundaries.

  \item[Kits are not templates or examples.] A Kit is not a collection of
    example runbooks.  It is a semantic layer: schemas, validators, a compiler,
    projections, and testing contracts.  An example runbook may demonstrate a
    Kit's vocabulary, but the Kit itself is the machinery that makes that
    vocabulary valid and executable.

  \item[Kits do not own execution.] The core runtime is the sole authority
    over step advancement, state management, governance enforcement, evidence
    capture, and trace writing.  A Kit may influence what the runtime executes
    (via lowering), but it never participates in how execution proceeds.

  \item[Kits do not replace the core schema.] The core \texttt{runbook/v1}
    schema remains the stable contract between authors, the parser, and the
    runtime.  Kit schemas are overlays that compile down to core; they do not
    bypass or replace it.

  \item[Kits are not cloud-native constructs.] Consistent with gert's
    local-first philosophy, Kits do not require network access, remote
    registries, or cloud services.  A Kit is a local artifact---a directory
    with a manifest, a compiler, and schemas---that runs wherever gert runs.
\end{description}

Kit discovery and distribution mechanisms are covered in Chapter~\ref{sec:kit-catalog}.
